% GENERATED FILE. Edit canonical lessons, then run npm run generate:books.
% canonical-source-hash: fnv1a64:240fec846ff09657
% canonical-lessons: SA-C116-sarasah, SA-C116-kapolah, SA-C116-cibukam, SA-C116-kurparah, SA-C116-manibandhah

\chapter{Crane, Cheek, Chin, Elbow, Wrist}
\label{ch:sa-crane-cheek-chin-elbow-wrist}
\hlchaptermodality{116}

\begin{chapteropening}
\textbf{By the end of this chapter:} \emph{I can say a crane, a cheek, the chin, an elbow and a wrist.}

Complete the last lesson of chapter 116: i can say a crane, a cheek, the chin, an elbow and a wrist.
\end{chapteropening}

\section[sārasaḥ]{\sk{सारसः} (sārasaḥ) — a crane}

\label{lesson:SA-C116-sarasah}



\begin{quote}
Before the new one: say the Sanskrit for a butterfly, then the Sanskrit for a rooster.
\end{quote}

\subsection*{You'll want to know: \sk{सारसः}}
\textbf{\sk{सारसः}} — \emph{sārasaḥ} — \textquotedblleft{}a crane\textquotedblright{}.

\begin{grammarlens}[title={What you've built}]
One more word for this chapter.
\end{grammarlens}

\subsection*{Your turn}
\begin{itemize}
\raggedright
  \item \emph{Say it:} \emph{sārasaḥ}
  \item \emph{Say it:} \emph{sārasaḥ}, once more
  \item \emph{Say it:} say \emph{sārasaḥ}
  \item \emph{Recall:} say the Sanskrit for a butterfly, then the Sanskrit for a rooster, then say \emph{sārasaḥ} again
\end{itemize}

\begin{tcolorbox}[breakable,colback=teal!4,colframe=teal!35!black,title={Before you move on}]
What does \sk{सारसः} mean? (\textquotedblleft{}A crane\textquotedblright{}.) Say it once more.
\end{tcolorbox}
\section[kapolaḥ]{\sk{कपोलः} (kapolaḥ) — a cheek}

\label{lesson:SA-C116-kapolah}



\begin{quote}
Before the new one: say the Sanskrit for a rooster, then the Sanskrit for a crane.
\end{quote}

\subsection*{You'll want to know: \sk{कपोलः}}
\textbf{\sk{कपोलः}} — \emph{kapolaḥ} — \textquotedblleft{}a cheek\textquotedblright{}.

\begin{grammarlens}[title={What you've built}]
One more word for this chapter.
\end{grammarlens}

\subsection*{Your turn}
\begin{itemize}
\raggedright
  \item \emph{Say it:} \emph{kapolaḥ}
  \item \emph{Say it:} \emph{kapolaḥ}, once more
  \item \emph{Say it:} say \emph{kapolaḥ}
  \item \emph{Recall:} say the Sanskrit for a rooster, then the Sanskrit for a crane, then say \emph{kapolaḥ} again
\end{itemize}

\begin{tcolorbox}[breakable,colback=teal!4,colframe=teal!35!black,title={Before you move on}]
What does \sk{कपोलः} mean? (\textquotedblleft{}A cheek\textquotedblright{}.) Say it once more.
\end{tcolorbox}
\section[cibukam]{\sk{चिबुकम्} (cibukam) — the chin}

\label{lesson:SA-C116-cibukam}



\begin{quote}
Before the new one: say the Sanskrit for a crane, then the Sanskrit for a cheek.
\end{quote}

\subsection*{You'll want to know: \sk{चिबुकम्}}
\textbf{\sk{चिबुकम्}} — \emph{cibukam} — \textquotedblleft{}the chin\textquotedblright{}.

\begin{grammarlens}[title={What you've built}]
One more word for this chapter.
\end{grammarlens}

\subsection*{Your turn}
\begin{itemize}
\raggedright
  \item \emph{Say it:} \emph{cibukam}
  \item \emph{Say it:} \emph{cibukam}, once more
  \item \emph{Say it:} say \emph{cibukam}
  \item \emph{Recall:} say the Sanskrit for a crane, then the Sanskrit for a cheek, then say \emph{cibukam} again
\end{itemize}

\begin{tcolorbox}[breakable,colback=teal!4,colframe=teal!35!black,title={Before you move on}]
What does \sk{चिबुकम्} mean? (\textquotedblleft{}The chin\textquotedblright{}.) Say it once more.
\end{tcolorbox}
\section[kūrparaḥ]{\sk{कूर्परः} (kūrparaḥ) — an elbow}

\label{lesson:SA-C116-kurparah}



\begin{quote}
Before the new one: say the Sanskrit for a cheek, then the Sanskrit for the chin.
\end{quote}

\subsection*{You'll want to know: \sk{कूर्परः}}
\textbf{\sk{कूर्परः}} — \emph{kūrparaḥ} — \textquotedblleft{}an elbow\textquotedblright{}.

\begin{grammarlens}[title={What you've built}]
One more word for this chapter.
\end{grammarlens}

\subsection*{Your turn}
\begin{itemize}
\raggedright
  \item \emph{Say it:} \emph{kūrparaḥ}
  \item \emph{Say it:} \emph{kūrparaḥ}, once more
  \item \emph{Say it:} say \emph{kūrparaḥ}
  \item \emph{Recall:} say the Sanskrit for a cheek, then the Sanskrit for the chin, then say \emph{kūrparaḥ} again
\end{itemize}

\begin{tcolorbox}[breakable,colback=teal!4,colframe=teal!35!black,title={Before you move on}]
What does \sk{कूर्परः} mean? (\textquotedblleft{}An elbow\textquotedblright{}.) Say it once more.
\end{tcolorbox}
\section[maṇibandhaḥ]{\sk{मणिबन्धः} (maṇibandhaḥ) — a wrist}

\label{lesson:SA-C116-manibandhah}



\begin{quote}
Before the new one: say the Sanskrit for the chin, then the Sanskrit for an elbow.
\end{quote}

\subsection*{You'll want to know: \sk{मणिबन्धः}}
\textbf{\sk{मणिबन्धः}} — \emph{maṇibandhaḥ} — \textquotedblleft{}a wrist\textquotedblright{}.

\begin{grammarlens}[title={What you've built}]
One more word for this chapter.
\end{grammarlens}

\subsection*{Your turn}
\begin{itemize}
\raggedright
  \item \emph{Say it:} \emph{maṇibandhaḥ}
  \item \emph{Say it:} \emph{maṇibandhaḥ}, once more
  \item \emph{Say it:} say \emph{maṇibandhaḥ}
  \item \emph{Recall:} say the Sanskrit for the chin, then the Sanskrit for an elbow, then say \emph{maṇibandhaḥ} again
\end{itemize}

\begin{tcolorbox}[breakable,colback=teal!4,colframe=teal!35!black,title={Before you move on}]
What does \sk{मणिबन्धः} mean? (\textquotedblleft{}A wrist\textquotedblright{}.) Say it once more.
\end{tcolorbox}
